\section{Background and Related Work}
\label{sec:background}

\subsection{The Evolution and Limitations of Blueprints}
\label{sec:bg-blueprints}

% - Origins: Scholze's Liquid Tensor Experiment, Massot's leanblueprint
% - Success: PFR (Tao 2023) — blueprint-first, distributed collaboration, color progress tracking
% - LeanArchitect (Zhu 2026): extract blueprint metadata from Lean code, eliminate LaTeX/Lean duplication, but direction is Lean → blueprint not blueprint → Lean
% - Common limitations: untyped edges, no network analysis, dual-track maintenance

\subsection{Network Science Perspectives on Mathematical Knowledge}
\label{sec:bg-network-science}

% - Mathlib network structure analysis (author's prior work): multi-layer dependency networks, community structure, centrality distributions, hierarchical features
% - Network science methods in other knowledge systems (citation network analysis, knowledge graph completion)
% - Key insight: topological features of mathematical knowledge networks can guide formalization strategy — high betweenness centrality nodes are bottlenecks, prioritizing their formalization unlocks the most downstream work

\subsection{Content-Addressing and Decentralized Formalization}
\label{sec:bg-content-addressing}

% - Stacks Project tag system: content-addressing pioneer, but flat data model
% - Content-addressing pipeline for formal mathematics (author's prior work)
% - IPFS/CID paradigm successes in package management

\subsection{Category Theory in Data Modeling}
\label{sec:bg-category-theory}

% - Spivak's Ologs (2012): most direct theoretical precursor
%   - Olog = category + natural language labels
%   - Objects are types, morphisms are functional relations (aspects)
%   - Commutative diagrams are facts, functors are mappings between knowledge bases
%   - Key contributions: (1) commutative diagrams express path equivalence; (2) functors enable precise KB translation; (3) Ologs convert to DB schemas
% - Spivak's Functorial Data Migration: functor-based data migration between DB schemas
% - Multi-model database categorical unification (Koupil et al. 2022, 2025)
% - Topos Institute: Double Categories for Knowledge Representation
% - Category Theory and Model-Driven Engineering (Diskin & Maibaum)
% - Common limitation: all remain at theoretical/formal level, no end-user interactive KM tools produced
% - Astrolabe's position: not theorem proving, not DB theory, but engineering Ologs' core insight (obj/mor as first-class citizens) — lower barrier (AI), relax constraints (typed relations instead of functions), add dimensions Ologs didn't address (content-addressing, network analysis, plugins, visualization)
